% ─────────────────────────────────────────────────────────────────────────────
%  Abstract page — Bachelor's Thesis (TFG)
%  Knowledge Distillation and Efficient Deployment of Language Models
%  in High-Performance Computing Environments
%  Marc Romera · BSC · UPC · 2026
%  Compile:  pdflatex abstract.tex
% ─────────────────────────────────────────────────────────────────────────────

\documentclass[11pt, a4paper]{article}

% ── Encoding & language ───────────────────────────────────────────────────────
\usepackage[utf8]{inputenc}
\usepackage[T1]{fontenc}
\usepackage[english]{babel}

% ── Fonts & typography ────────────────────────────────────────────────────────
\usepackage{lmodern}
\usepackage{microtype}

% ── Layout ───────────────────────────────────────────────────────────────────
\usepackage[
  a4paper,
  top    = 2.0cm,
  bottom = 1.8cm,
  left   = 3.0cm,
  right  = 3.0cm
]{geometry}

% ── Spacing ───────────────────────────────────────────────────────────────────
\usepackage{setspace}
\usepackage{parskip}
\setlength{\parindent}{0pt}
\setlength{\parskip}{5pt}

% ── Color ────────────────────────────────────────────────────────────────────
\usepackage[dvipsnames]{xcolor}
\definecolor{thesisblue}{HTML}{1B3A6B}
\definecolor{darkgray}{HTML}{333333}
\definecolor{medgray}{HTML}{888888}

% ── Rules & misc ─────────────────────────────────────────────────────────────
\usepackage[hidelinks]{hyperref}

% ─────────────────────────────────────────────────────────────────────────────
\begin{document}
\pagestyle{empty}

% ════════════════════════════════════════════════════════════════════════════
%  HEADER BLOCK
% ════════════════════════════════════════════════════════════════════════════

{\color{thesisblue}\rule{\linewidth}{2.4pt}}

\vspace{0.4cm}

% ── Thesis type label ────────────────────────────────────────────────────────
\begin{center}
  {\small\color{medgray}\textsc{Bachelor's Thesis in Computer Science}}
\end{center}

\vspace{0.25cm}

% ── Title ────────────────────────────────────────────────────────────────────
\begin{center}
  {\color{thesisblue}
   \fontsize{14.5}{18}\selectfont\bfseries
   Knowledge Distillation and Efficient Deployment of Language Models\\[4pt]
   in High-Performance Computing Environments}
\end{center}

\vspace{0.25cm}

% ── Author / institution / supervisor ────────────────────────────────────────
\begin{center}
  {\small\color{darkgray}
    \textbf{Marc Romera}\\[3pt]
    Barcelona Supercomputing Center (BSC) \quad
    Universitat Polit\`ecnica de Catalunya (UPC)\\[3pt]
    Barcelona, March 2026}
\end{center}

\vspace{0.3cm}
{\color{medgray}\rule{\linewidth}{0.5pt}}

% ════════════════════════════════════════════════════════════════════════════
%  ABSTRACT
% ════════════════════════════════════════════════════════════════════════════

\vspace{0.3cm}

\begin{center}
  {\color{thesisblue}\large\bfseries Abstract}
\end{center}

\vspace{0.15cm}

\begin{singlespace}
\color{darkgray}\small

The computational cost of serving large language models at scale represents
a significant barrier to their wider adoption.
A central motivation of this thesis is that request complexity varies
considerably in practice, and many queries can be adequately resolved by
smaller, more efficient models.
This suggests an opportunity to design inference systems that dynamically
assign each request to the smallest model capable of meeting a defined
quality standard, reducing average serving cost while aiming to preserve
output quality.

This thesis designs, implements, and evaluates such a system, built around
a hierarchy of instruction tuned models from the Qwen2.5 family (14B, 7B,
and 1.5B parameters) and deployed on MareNostrum~5 at the Barcelona
Supercomputing Center using SLURM-orchestrated NVIDIA H100 GPUs and vLLM
with OpenAI compatible endpoints.
Three serving policies are studied under a common experimental framework.
A baseline policy routes all requests to the most capable model.
A forced cascade attempts each request with the smallest model and escalates
progressively until a quality criterion is satisfied.
A confidence based cascade uses output log-probability signals to estimate
which minimum model is likely to meet the quality threshold and routes the
request accordingly, reducing unnecessary escalation.
In parallel, knowledge distillation via supervised fine tuning with LoRA
adapters is applied to narrow the capability gap between the student models
and the teacher, using teacher-generated reasoning traces on mathematical
tasks as training supervision.

Quality is assessed on GSM8K and MATH-500; system efficiency is
characterised through throughput and latency measurements conducted under
reproducible, controlled conditions on MareNostrum~5.
Together, the implemented system and evaluation framework provide a
reproducible basis for studying distillation and adaptive routing strategies
across the quality-efficiency trade-off in large language model serving at
HPC scale.

\end{singlespace}

% ════════════════════════════════════════════════════════════════════════════
%  KEYWORDS
% ════════════════════════════════════════════════════════════════════════════

\vspace{0.3cm}
{\color{medgray}\rule{\linewidth}{0.5pt}}
\vspace{0.2cm}

{\small\color{darkgray}
  \textbf{\color{thesisblue}Keywords}\enspace
  \textit{%
    knowledge distillation,
    large language models,
    high-performance computing,
    efficient inference,
    vLLM,
    LoRA,
    throughput,
    latency,
    MareNostrum~5,
    quality--efficiency tradeoff%
  }
}

\vspace{0.3cm}
{\color{thesisblue}\rule{\linewidth}{2.4pt}}

\end{document}
